% !PRINTED 11
\weberpage{11}

%<<CONT
\gesperrt{Wenn $Q$ ein Theiler von $P$ ist, der mit seinen sämmtlichen conjugirten
Theilern identisch ist, so heisst $Q$ ein Normaltheiler von
$P$}\footnote{Auch ausgezeichnete oder invariante Untergruppe genannt.}.

Die aus dem einzigen Element 1 gebildete Gruppe ist ein Normaltheiler von
jeder Gruppe. Wir erhalten ferner einen Normaltheiler in dem grössten
gemeinschaftlichen Theiler $R$ der sämmtlichen mit irgend einem Theiler
$Q$ von $P$ conjugirten Theiler.

Denn es ist schon oben bewiesen, dass $R$ als der Durchschnitt mehrerer
Theiler von $P$ eine Gruppe ist.

Sind nun
\weq{1}{Q,\ Q',\ Q'' \ldots}
die zu $Q$ conjugirten Theiler von $P$ und $b$ irgend ein Element von $P$,
dann ist das System der Gruppen
\weq{2}{b^{-1}Qb,\ b^{-1}Q'b,\ b^{-1}Q''b \ldots}
von dem System (1) nicht verschieden. Wenn aber $R$ der Durchschnitt der
Gruppen (1) ist, so ist $b^{-1}Rb$ der Durchschnitt von (2), und folglich
ist $R$ mit $b^{-1}Rb$ identisch, d.\,h. $R$ ist ein Normaltheiler von $P$.

Ist $N$ ein Normaltheiler irgend einer Gruppe $P$ und $b$ ein beliebiges
Element in $P$, so ist
\weq{3}{b^{-1}Nb = N \quad \text{oder} \quad Nb = bN.}

Ist $P$ vom Grade $n$, $N$ vom Grade $\nu$ und vom Index $\mu$, also
$n = \mu\nu$, so können wir die $\mu$ Elemente $1, b_1 \ldots b_{\mu-1}$ so
wählen, dass die Zerlegung von $P$ in die Nebengruppen
\[
\begin{aligned}
P &= N + Nb_1 + Nb_2 + \cdots + Nb_{\mu-1} \\
  &= N + b_1 N + b_2 N + \cdots + b_{\mu-1} N
\end{aligned}
\]
ergiebt. Es ist also
\weq{4}{N,\ N_1 = Nb_1 = b_1 N,\ N_2 = Nb_2 = b_2 N \ldots,\quad N_{\mu-1} = Nb_{\mu-1} = b_{\mu-1} N}
das System der Nebengruppen.

Ist $a$ irgend ein Element in $N$, so ist $Na$ mit $N$ identisch, also ist
auch $Nab$ mit $Nb$ identisch, und da jedes Element $c$ von $P$ in $N$
oder in einer seiner Nebengruppen vorkommt, also in der Form $ab_k$
enthalten ist, so muss $Nc = cN$ mit einem der Systeme (4)
übereinstimmen.
